\subsection{Rietsch's Lie-theoretical mirror}

All varieties are over \(\CC\). We fix the notation from Rietsch's Lie-theoretic mirror construction. Let \(G\) be a connected simply connected semisimple algebraic group. Fix opposite Borel subgroups \(B_+\) and \(B_-\), with unipotent radicals \(U_+\) and \(U_-\), and put
\[
T=B_+\cap B_-.
\]
The Weyl group is
\[
W=N_G(T)/T,
\]
with longest element \(w_0\). For \(w\in W\), we choose a representative in \(G\), denoted by \(\dot w\).

Let \(\mathfrak g,\mathfrak b_\pm,\mathfrak u_\pm,\mathfrak h\) be the Lie algebras of \(G,B_\pm,U_\pm,T\). Let \(\Delta_+\) be the set of positive roots determined by \(B_+\). We write \(I\) for the index set of simple roots \(\{\alpha_i:i\in I\}\). For \(i\in I\), choose Chevalley generators
\[
e_i\in\mathfrak g_{\alpha_i},
\qquad
f_i\in\mathfrak g_{-\alpha_i}.
\]
They define one parameter subgroups
\[
x_i(a)=\exp(ae_i),
\qquad
y_i(a)=\exp(af_i).
\]
We take
\[
\dot s_i=x_i(1)y_i(-1)x_i(1)
\]
as the representative of the simple reflection \(s_i\).

Let \(P\supset B_-\) be a parabolic subgroup. Define
\[
I_P=\{i\in I:\dot s_i\in P\},
\qquad
I^P=I\setminus I_P.
\]
The parabolic Weyl subgroup is
\[
W_P=\langle s_i:i\in I_P\rangle,
\]
with longest element \(w_P\). We also use
\[
W^P=\{w\in W:\ell(ws_i)>\ell(w)\text{ for all }i\in I_P\}
\]
for the minimal length representatives of \(W/W_P\). If
\[
I^P=\{n_1,\ldots,n_k\},
\qquad
n_1<\cdots<n_k,
\]
then \(k=\dim H^2(G/P,\CC)\).

Let \(G^\vee\) be the Langlands dual group. Since \(G\) is simply connected, \(G^\vee\) is adjoint. We use the notation \(B_\pm^\vee,U_\pm^\vee,T^\vee,\mathfrak g^\vee\), and Chevalley generators \(e_i^\vee,f_i^\vee\). The Weyl group is again \(W\), and we keep the dot notation \(\dot w\) for representatives in \(G^\vee\) when no confusion can occur.

For \(v\leq w\) in the Bruhat order, the open Richardson variety is the intersection of a Schubert variety and an opposite Schubert variety in \(G^\vee/B_-^\vee\):
\[
\mathcal R_{v,w}^\vee
=
\left(B_+^\vee\dot vB_-^\vee\cap B_-^\vee\dot wB_-^\vee\right)/B_-^\vee
\subset G^\vee/B_-^\vee .
\]
It is smooth and irreducible of dimension \(\ell(w)-\ell(v)\). The open Richardson variety relevant to \(G/P\) is \(\mathcal R_{w_P,w_0}^\vee\). Let \(\pi:G^\vee/B_-^\vee\to G^\vee/P^\vee\) be the natural projection. Its image is a projected Richardson stratum in the sense of Knutson--Lam--Speyer \cite{KLS}. We denote the reduced complement of this image by
\[
D_{\mathrm{Lie}}^\vee=G^\vee/P^\vee\setminus \pi(\mathcal R_{w_P,w_0}^\vee).
\]
The divisor \(D_{\mathrm{Lie}}^\vee\) is the Rietsch boundary, namely the boundary selected by the Lie-theoretical mirror.

The parameters come from the small quantum cohomology ring of \(G/P\), which is a deformation of \(H^*(G/P)\) with
\[
k=\dim H^2(G/P,\CC)
\]
parameters \(q_1,\ldots,q_k\). These parameters correspond to a basis of curve classes dual to the Schubert divisor classes. In the above Lie-theoretic notation, they are indexed by
\[
I^P=\{n_1,\ldots,n_k\}.
\]

Rietsch identifies the nonzero quantum parameter space with the torus \((T^\vee)^{W_P}\). More precisely,
\[
(T^\vee)^{W_P}=\{t\in T^\vee:\alpha_i^\vee(t)=1\ \text{for all }i\in I_P\},
\]
and the coordinates \(q_j\) are identified with the characters \(\alpha_{n_j}^\vee\) on this torus. Thus working over \((T^\vee)^{W_P}\) is the same as inverting the quantum parameters.

%Rietsch's work gives a uniform Lie theoretic mirror construction for general flag varieties \(G/P\). It was motivated by the earlier mirror models of Givental for toric Fano varieties and type \(A\) flag varieties, and by Peterson's presentation of quantum cohomology. In the full equivariant setting, the mirror data also include a holomorphic volume form and a multi-valued function. For our purpose, the important part is the family together with the superpotential, because this is what determines the open Richardson domain and the compactification boundary.

For fixed \(t\in (T^\vee)^{W_P}\), we have 
\[
Z_P^t=B_-^\vee\cap U_+^\vee\,t\dot w_P\dot w_0^{-1}U_+^\vee .
\]
These fibres form a family
\[
Z_P=
\{(t,b)\in (T^\vee)^{W_P}\times B_-^\vee:
b\in U_+^\vee\,t\dot w_P\dot w_0^{-1}U_+^\vee\}
\]
over \((T^\vee)^{W_P}\) by projection to the first factor. 

Rietsch proves that \(Z_P^t\) is smooth of dimension \(\dim G/P\), and the map
\[
Z_P^t\longrightarrow \mathcal R_{w_P,w_0}^\vee,\qquad
b\longmapsto b\dot w_0B_-^\vee/B_-^\vee
\]
is an isomorphism. Hence each fibre has a natural compactification by \(G^\vee/P^\vee\), with boundary \(D_{\mathrm{Lie}}^\vee\).

Define \(F=\sum_{i\in I}(e_i^\vee)^*\in(\mathfrak g^\vee)^*\), where \(e_i^\vee\) are the Chevalley generators fixed above. Let \(\rho\) denote the sum of the fundamental coweights. By definition, each \(b\in Z_P^t\) has a factorization
\[
b=u_1t\dot w_P\dot w_0^{-1}u_2^{-1},
\qquad u_1,u_2\in U_+^\vee,
\]
and Rietsch's superpotential \cite{Rietsch} is
\[
\mathcal F_P(t,b)=F(u_2\cdot\rho)-F(u_1\cdot\rho).
\]
This expression is equivalent to the representation-theoretic definition and gives a regular function on \(Z_P\).

The Rietsch mirror of \(G/P\) is the Landau--Ginzburg model \((Z_P,\mathcal F_P)\to (T^\vee)^{W_P}\). After fixing the quantum parameters \(t\), we can also write it as
\[
\left(X_P^\vee,\mathcal W_t\right)
=
\left(G^\vee/P^\vee\setminus D_{\mathrm{Lie}}^\vee,\ \mathcal F_P|_{Z_P^t}\right),
\]
where \(Z_P^t\simeq X_P^\vee\) through \(b\mapsto b\dot w_0B_-^\vee\). The critical locus of \(\mathcal F_P\) along the fibres recovers Peterson's presentation of the quantum cohomology of \(G/P\) with the quantum parameters inverted \cite{Peterson,LS}.

For the type \(C_n\) model considered below, the original homogeneous space is
\[
Sp_{2n}/P_1\simeq \PP^{2n-1}.
\]
The Rietsch mirror for this homogeneous space is defined on the dual odd quadric
\[
SO_{2n+1}/P_1^\vee\simeq Q^{2n-1}.
\]
The fibre \(Z_P^t\) is identified with the open Richardson variety \(\mathcal R_{w_P,w_0}^\vee\), and compactifying this open variety inside \(G^\vee/P_1^\vee\) produces the dual-side boundary
\[
D_{\mathrm{Lie}}^\vee
=
G^\vee/P_1^\vee\setminus \pi(\mathcal R_{w_P,w_0}^\vee).
\]
Following the philosophy of the Gross--Siebert program, the mirror map should be viewed as a relation between log Calabi--Yau pairs, not only between the ambient spaces. Thus, on the original projective space \(\PP^N\), we also use the open Richardson variety attached to the type \(C_n\) homogeneous model \(Sp_{2n}/P_1\). Its complement is the type \(C_n\) boundary \(D_C\subset\PP^N\) computed below. This construction follows Rietsch's Richardson-stratum viewpoint and the notation used by Pech, Rietsch, and Williams~\cite{Rietsch,PRW}.
